\chapter{研究背景}

数据治理研究逐步从单一的数据质量检查扩展到数据标准、职责分工、过程控制与共享机制。
相关研究认为，治理规则需要与实际业务过程结合，才能形成可执行、可复核的管理闭环
\yibincite{zhang2024}。

\section{国内研究进展}

国内研究重视数据标准、目录管理和质量评价，常以制度规范与信息系统协同为主要路径。

\section{国外研究进展}

国外研究较早关注数据所有权、元数据管理和数据生命周期，并形成了较系统的方法框架
\yibincite{wang2023}。

\chapter{研究评述}

现有成果提供了较完整的概念和方法基础，但面向具体公共服务流程的字段级规则、证据记录
和持续评价机制仍需细化。后续研究应加强治理要求与实际处理步骤之间的映射。
